\section{Requirements and System Design}

This chapter translates the research gap into a design that can be implemented and evaluated. It first defines the product position and release requirements, then compares rejected alternatives, traces the decision lifecycle, specifies the evidence and rule-pack contracts, and maps each design invariant to an acceptance source. This order follows the central dependency of the system: a user-facing claim is valid only when the evidence, rule, and state transitions that produced it remain reconstructable.

\subsection{Design Objectives and Requirements}

Privacy Lens is designed as a local review instrument. Its output is a structured prompt for investigation, not a legal decision and not a security-enforcement action. The system accepts bounded permission summaries, validates their shape and temporal relationships, applies the active regulation pack, stores the resulting findings locally, and presents the result with its provenance and limitations.

Table~\ref{tab:r11-requirements} defines the requirements used for implementation and evaluation. The table is intentionally concise; later subsections explain why each requirement is necessary and which failure it prevents.

\begin{table}[H]
\centering
\small
\caption{the final delivery requirements and acceptance evidence}
\label{tab:r11-requirements}
\begin{tabular}{p{0.13\textwidth}p{0.50\textwidth}p{0.27\textwidth}}
\toprule
ID & Requirement & Evidence \\
\midrule
F1 & Evaluate supported permission summaries deterministically. & Seeded and boundary suites \\
F2 & Preserve source, time window, pack, rule, and missing context. & Type and repository tests \\
F3 & Reject unsafe identifiers, numbers, sources, timestamps, windows, and contexts. & Adversarial tests \\
F4 & Switch only among registered, versioned regulation packs. & Pack registry tests and Settings UI \\
F5 & Separate synthetic demonstrations from device-visible evidence. & Source labels and simulator tests \\
U1 & Provide overview, findings, and settings through stable navigation. & Physical screenshots and UI trees \\
U2 & Communicate severity without colour alone and disclose evidence limits. & Source inspection and physical rendering \\
P1 & Keep findings local, disable Android backup, and request no sensitive runtime permission. & Manifest and installed package inspection \\
D1 & Build a versioned release APK and AAB with a stable package identity. & Release build and device install \\
\bottomrule
\end{tabular}
\end{table}

\subsection{Design Alternatives and Threat Model}

The model includes malformed bridge values, forged simulator provenance, unsafe numbers, replayed windows, overlapping summaries, unbounded history, misleading empty results, and user over-reliance on authoritative language. It does not defend against a compromised operating system, a maliciously modified application binary, falsified system telemetry, or a legally incorrect pack supplied by a trusted maintainer. Those risks require platform integrity, signing governance, and independent legal review.

\subsubsection{Selection Rationale}

Several alternative product positions were considered. The first was a binary compliance scanner that would display compliant or non-compliant for every application. This was rejected because permission frequency cannot establish purpose, necessity, lawful basis, controller obligations, or applicable jurisdiction. A binary label would hide missing evidence and encourage users to treat a technical heuristic as a legal decision. The adopted three-state model retains a needs-review state, an evidence-gap state, and a no-technical-concern state whose wording is explicitly limited.

The second alternative was a single risk score. A score appears compact and supports sorting, but it combines unlike uncertainties. A high aggregate count, missing purpose, synthetic source, and consent contradiction do not have a common natural unit. Weighting them would introduce assumptions that are difficult to justify and easy to overlook. Privacy Lens instead stores contributing signals and renders a reasoned state. A future prioritisation score could be added for professional triage only if every contribution remains inspectable and calibration evidence exists.

The third alternative was a learned classifier trained on application labels or permission histories. Machine learning could model interactions not captured by fixed thresholds, yet the project lacked an independently labelled, representative, legally meaningful data set. Training on store categories, malware labels, or developer declarations would answer different questions. A deterministic rule engine was selected because it makes boundary behaviour, missing inputs, and change history reviewable. This choice does not claim that rules are universally superior; it matches the evidence available for the thesis.

The fourth alternative was to embed GDPR language in interface components. This is initially convenient because each screen can display the relevant article directly. It creates duplication, complicates review, and makes regional extension risky: changing jurisdiction would require editing presentation and perhaps persistence logic throughout the application. The pack contract centralises legally motivated content and forces the engine to record which pack generated a finding. Generic screens display pack-provided labels through validated fields.

The fifth alternative was to choose a region automatically from device locale or location. Locale does not reliably determine applicable law, organisational establishment, data-subject status, or processing context. Automatic selection could therefore present the wrong framework with unwarranted confidence. The adopted design requires explicit user selection from registered packs and shows the selected jurisdiction and caveat. It does not collect location to decide which privacy law to discuss.

The sixth alternative was a cloud account with central evidence storage. Centralisation could simplify synchronisation, remote dashboards, and model updates, but it would also create a new repository of sensitive behavioural metadata and require authentication, server security, retention, international transfer analysis, and deletion workflows. The current product is local first, uses no account, and offers local deletion. This choice reduces deployment features but aligns the application's conduct with its privacy objective.

The seventh alternative was to ingest unrestricted free-form pack files. This would maximise extensibility but expose users to misleading sources, malformed thresholds, pack confusion, and potentially executable content. The present registry admits bundled immutable definitions. Future signed declarative distribution may be appropriate, but only after schema, publisher, rollback, and review controls exist.

The eighth alternative was to treat absence of observations as zero activity. This would produce a cleaner dashboard, but it would conflate a real zero with unavailable acquisition. The selected model requires a capability-aware source and reports an evidence gap when coverage is missing. This is less reassuring and more accurate.

The ninth alternative concerned notifications. Emitting an alert for every evaluation would demonstrate activity but create repetition and fatigue. Repository logic instead considers whether a finding is new, has changed state, or has escalated. Silent storage remains possible for unchanged or low-priority outcomes. The design can later support controlled comparison of notification policies without changing decision semantics.

The final alternative concerned storage of raw observations. Keeping a complete event history would improve forensic detail but increase sensitivity and resource cost. The candidate stores bounded findings and aggregate history needed for current temporal reasoning. Event-level acquisition is deferred until an authority model, minimisation policy, and retention design can be evaluated. This is a deliberate privacy-performance trade-off rather than an accidental omission.

\subsection{Architecture and Decision Lifecycle}

The selected architecture separates acquisition, admission, temporal preparation, evaluation, persistence, projection, and presentation. Rather than allowing each screen to interpret raw Android data independently, every stage receives a bounded input and produces a state whose meaning is explicit. The following lifecycle, failure, minimisation, and release-gate decisions define how this separation behaves under both normal and adverse conditions.

Figure~\ref{fig:system-architecture} shows the single decision path and the two governance inputs that constrain it. The regulation pack does not receive UI state, and the interface does not bypass the finding record to reinterpret raw evidence. This separation makes source identity and pack identity available at every downstream review point.

\begin{figure}[H]
\centering
\resizebox{0.94\textwidth}{!}{%
\begin{tikzpicture}[
  node distance=0.43cm,
  stage/.style={draw, rounded corners, minimum width=6.7cm, minimum height=0.72cm, align=center, font=\small},
  side/.style={draw, dashed, rounded corners, text width=3.5cm, align=center, font=\small},
  arrow/.style={-latex, thick}
]
\node[stage, fill=blue!7] (acq) {1. Acquisition: native bridge, imported record, or labelled simulator};
\node[stage, fill=blue!7, below=of acq] (admit) {2. Admission: owned identifiers, safe numbers, source, time, and context};
\node[stage, fill=orange!9, below=of admit] (temporal) {3. Temporal preparation: replay isolation and completed non-overlapping history};
\node[stage, fill=orange!9, below=of temporal] (evaluate) {4. Rule evaluation: threshold and context classification under one pack};
\node[stage, fill=green!9, below=of evaluate] (finding) {5. Finding construction: evidence, rationale, missing facts, caveat, and pack version};
\node[stage, fill=green!9, below=of finding] (store) {6. Local persistence and dashboard projection};
\node[stage, fill=purple!8, below=of store] (present) {7. Presentation: Overview, Findings, Settings, and human next action};
\foreach \a/\b in {acq/admit,admit/temporal,temporal/evaluate,evaluate/finding,finding/store,store/present}{\draw[arrow] (\a) -- (\b);}
\node[side, right=1.0cm of admit] (econtract) {Evidence contract\\provenance and validity boundary};
\node[side, right=1.0cm of evaluate] (pack) {Regulation-pack contract\\authority, version, rules, caveat};
\draw[dashed,-latex] (econtract) -- (admit);
\draw[dashed,-latex] (pack) -- (evaluate);
\end{tikzpicture}
}
\caption{Evidence-bounded system architecture and decision path}
\label{fig:system-architecture}
\end{figure}

\subsubsection{End-to-End Decision Lifecycle}

The lifecycle begins when the user initiates review, a scheduled worker requests work, or the controlled simulator generates an example. The acquisition component returns candidate records together with a source identity. It must not return a native zero merely because a platform capability is unavailable. Capability failure and an observed zero are different states and lead to different user decisions.

Each candidate crosses the admission boundary as unknown data. The validator checks the package identifier, permission key, count, window, timestamps, source, and optional context without relying on compile-time types. Rejection is structured. The record does not enter temporal history and cannot contribute to a reassuring status. The caller can present the failure as an evidence problem and retain enough diagnostic detail to correct the producer without exposing raw sensitive content unnecessarily.

An admitted record is placed in a key space defined by package and permission. Exact-window replay replaces the prior aggregate. Independent completed windows may be retained for cross-window analysis; overlapping aggregates are not summed. Simulator data remains excluded from accumulated cross-window claims. Retention and a per-key cap bound state. These operations prepare evidence before jurisdictional rules are applied.

The engine receives the selected immutable pack explicitly. It locates the supported rule, calculates daily excess, evaluates burst information where timestamps are available, examines compatible history, and inspects any admitted processing context. The result includes the observed and threshold values, signal types, status, communication priority, rationale, legal or research reference, missing-context questions, and pack identity. The calculation never mutates the pack.

The repository compares the result with the stored finding for the same logical subject. It distinguishes an unchanged state, a new concern, resolution or reversal, and risk escalation. This comparison governs notification eligibility. It also writes a bounded slice under the versioned repository key. A finding remains attached to its original pack even if Settings later selects another pack.

Projection turns repository records into interface-facing models. It calculates overview totals by state, applies Findings filters, formats source and temporal labels, and exposes actions. Projection should not infer new legal semantics. If a required field is absent, it retains an unknown or evidence-gap presentation. This makes projection testable independently of React components.

Presentation then orders information according to user need. The overview states reach and limitation before the main action. Findings allow scanning by state and inspection of reasons. Settings places pack choice, caveat, simulator, privacy information, and deletion controls in separate groups. Accessibility labels name the action and state rather than depending on icon or colour. Destructive action requires deliberate confirmation where implemented.

The lifecycle ends with user action or explicit non-action. A user may seek a privacy notice, ask a developer for purpose and retention information, repeat review with authorised evidence, or decide that the technical signal does not require immediate attention. Privacy Lens does not contact the reviewed application, change its permissions, or submit a complaint automatically. Keeping the final decision outside the engine preserves human accountability.

Deletion is also part of the lifecycle. The user can clear local findings, and uninstall or platform data clearing removes the application data within Android's boundary because backup is disabled. Export or future synchronisation would create additional copies and must define deletion separately. The current lifecycle intentionally ends at local storage.

\subsubsection{Failure-State Design}

Failure states were designed as product states rather than exception messages. Acquisition unavailable means the current deployment cannot obtain suitable records; it is not displayed as a completed scan. Invalid evidence means a producer supplied data that fails the contract; it cannot be silently repaired into a normal case. Unknown pack means evaluation is blocked until a registered definition is selected; it does not fall back to the EU pack merely because that pack is bundled.

Storage failure also requires conservative behaviour. If persisted data cannot be parsed or migrated safely, the interface should not mix it with current findings. Recoverable records may be retained after field-level validation, while unrecoverable content should be quarantined or omitted with a visible explanation. Clearing data remains available so the user can recover control. Production logging should report the failure class without printing evidence payloads.

Work scheduling failure is distinct from review failure. WorkManager can defer periodic work because of platform policy, Doze, reboot, or vendor restrictions. The interface should present the last completed review time and source capability rather than imply continuous monitoring. A registered request demonstrates scheduling intent, not on-time execution. Future long-run tests should measure delay and recovery before stronger wording is used.

Partial evaluation is another possibility. Some records may be valid while others are rejected. Treating the batch as wholly clean would hide rejection; treating it as wholly failed would discard useful evidence. A suitable summary reports accepted findings together with the count and reasons for excluded records. The current bounded flows establish structured rejection at record level; richer batch disclosure remains an extension point.

Pack-content failure has two forms. Schema failure is machine detectable and should block admission. Substantive legal error may only be identified by expert review or later correction. Pack provenance, versions, official sources, and historical identity allow the maintainer to issue a successor without pretending earlier findings never existed. The interface can mark a deprecated pack and recommend re-evaluation while preserving the original record.

Network failure is intentionally limited because core review and bundled packs operate locally. Privacy-policy and official-source links may be unavailable, but the application should retain the relevant title and caveat. If remote pack delivery or synchronisation is later added, offline operation and last-known-good behaviour become acceptance requirements rather than optional resilience features.

Visual failure includes clipped localisation, unreadable contrast, focus loss, hidden caveats, and controls obscured by system insets or keyboards. These failures may not appear in logic tests. Rendered screenshot matrices, hierarchy inspection, screen-reader use, font scaling, and device variation are therefore part of the planned evidence. The one-device acceptance check establishes only the candidate baseline.

Finally, misunderstanding is treated as a failure even when the software behaves as coded. If users interpret no technical concern as GDPR compliance, or mistake simulator records for device evidence, the product has failed its communication objective. Participant comprehension and reliance tests are needed to detect that class. The design reduces the risk through terminology and provenance, but does not declare it eliminated.

\subsubsection{Data-Minimisation Decisions}

The application needs enough evidence to explain a review signal, yet every retained field has privacy cost. Package identifier, permission family, count, time window, source, pack identity, and decision reason form the current minimum for reconstructing an aggregate finding. Optional context is admitted only when supplied for the review purpose. The application does not require account identity, contact lists, precise location content, microphone content, or advertising identifiers.

Counts are stored instead of raw sensor values. A location access finding does not require coordinates, and a microphone access finding does not require audio. This distinction is central: the system reviews evidence that an operation occurred, not the personal data produced by the operation. A future connector should preserve that separation unless a separately approved research question requires content.

Time resolution is also minimised. Aggregate start and end values support window validity, replay, and bounded temporal analysis. Timestamp arrays are accepted for controlled burst evidence, but they should not be collected by default merely because the schema can represent them. The source capability and purpose determine whether additional resolution is justified.

Local-first storage avoids creating a server-side behavioural database. The trade-off is that the user or controlled deployment is responsible for device access and loss, and professional collaboration features are limited. The design disables Android backup so evidence is not restored unexpectedly to another device through ordinary platform backup. This choice should be re-evaluated if enterprise-managed backup becomes a stated requirement.

Bounded history limits denial-of-service and retention exposure. The cap is not a substitute for purpose-based deletion, but it prevents indefinite growth under current aggregate use. Repository slices and clear controls make the lifecycle understandable. Future event-level storage would require stronger indexing, encryption assessment, retention schedules, and selective deletion.

Notifications contain only the minimum text needed to indicate a changed review state. Detailed evidence remains inside the application. Lock-screen visibility can still expose the existence of a concern, so future production settings should consider notification privacy and user preference. The current evaluation does not claim that every device lock-screen configuration was tested.

External links are used for official materials and privacy documentation instead of embedding broad web content. The application does not send the finding to those sites. Users should be informed that opening a link transfers them to another controller and may reveal ordinary browser metadata. Remote analytics or link tracking would contradict the minimisation posture unless separately justified.

Data minimisation finally applies to evaluation artefacts. Screenshots, UI dumps, log excerpts, and device metadata should avoid unrelated notifications, account names, and installed-application lists. Synthetic demonstrations are preferred for public evidence. Physical records used for thesis acceptance remain scoped to the test application and bounded device session.

\subsubsection{Requirement Priorities and Release Gates}

Requirements were prioritised according to the harm caused by failure. Evidence admission, provenance, non-verdict language, source separation, replay safety, pack identity, and deletion are release blockers because failure could create misleading privacy conclusions or remove user control. Core navigation, readable hierarchy, package identity, release builds, and absence of sensitive runtime permissions are also blockers for the initial store candidate.

Important but non-blocking research extensions include multi-OEM acquisition measurement, long-duration WorkManager observation, full screen-reader matrices, participant studies, production signing, and independent legal review. They are non-blocking only for the bounded engineering candidate. They become blockers for stronger claims such as public production readiness, accessibility conformance, or validated GDPR guidance.

Optional features include cloud synchronisation, accounts, remote pack updates, application blocking, and learned ranking. These features do not repair current evidence gaps automatically and would introduce substantial security, privacy, or governance work. They should be considered only after the core assurance chain is stable.

The release gate follows dependency order. A source-level failure blocks package evaluation because the binary would not represent an accepted program. A package or manifest failure blocks device acceptance. A device crash or broken primary journey blocks visual promotion. A missing legal or usability study does not turn an engineering check red; it keeps the associated substantive claim explicitly unverified.

This priority model avoids two extremes. It does not delay all prototype learning until every production activity is complete, and it does not allow a successful demonstration to erase unresolved authority or legal questions. Each gate authorises a named next activity. The the final delivery candidate is designed to authorise controlled store-submission preparation, internal-track testing, and participant evaluation.

\subsection{Stakeholders and Functional Requirements}

The same evidence contract serves different users, but their decisions and information needs are not identical. Requirements are therefore derived from explicit scenarios before they are divided into functional and non-functional categories.

\subsubsection{Stakeholders and Usage Scenarios}

The primary user is a person who wants to understand available privacy evidence without becoming an Android or GDPR specialist. This user needs a short account of what was observed, why it deserves attention, what is not known, and what action is reasonable. A second stakeholder is a developer or product owner using the application during internal review. This stakeholder needs reproducible rule identifiers, pack versions, timestamps, and evidence-source labels. A third stakeholder is a privacy or legal reviewer who may use the finding as an entry point into a broader assessment. This stakeholder needs the official source, missing-context questions, and an explicit statement that the application has not reached a legal conclusion.

The first usage scenario is an ordinary review with no prior data. The overview explains the scope before offering a review action. If the bridge returns no admissible observations, the interface reports an evidence gap rather than a clean bill of health. The second scenario is a populated review. Findings are grouped by status and can be inspected individually. The third scenario is learning through synthetic data. The user runs a controlled demonstration, which creates a varied but deterministic set of records marked as synthetic. The fourth scenario is changing legal perspective. The user selects another registered pack, sees its jurisdiction and caveat, and performs a new review without relabelling historical findings. The fifth scenario is withdrawal from the tool itself: the user clears local findings or uninstalls the application.

These scenarios exclude unsupported uses. Privacy Lens is not intended to block another application, intercept network traffic, prove that a sensor callback stopped, obtain evidence by evading platform restrictions, or submit a legal complaint automatically. It does not select applicable law from device location. Defining exclusions is important because a polished interface could otherwise suggest a broader product than the evidence pipeline implements.

\subsubsection{Functional Decomposition}

The review workflow is decomposed into acquisition, admission, temporal preparation, evaluation, persistence, projection, and presentation. Acquisition produces a candidate summary and a capability statement. Admission checks whether the record can be trusted as a well-formed input. Temporal preparation determines whether history can be combined without double counting. Evaluation applies one immutable pack. Persistence writes findings and migration metadata. Projection converts stored findings into dashboard counts and sections. Presentation renders this projection with actions and caveats.

This decomposition allows each failure to be represented at the correct layer. If the bridge is unavailable, the application should not create a zero-count native record. If a timestamp is unsafe, the engine should not try to repair it silently. If two aggregate windows overlap, the repository should not sum them. If a pack identifier is unknown, the registry should not fall back to whichever pack appears first. If storage migration fails, the interface should avoid presenting stale records as current evidence. The design favours explicit absence over invented continuity.

The simulator follows the same downstream path after acquisition. It is not a separate UI mock that bypasses validation. This is important because a demonstration that constructs presentation objects directly would not test the engine, repository, or projection. Synthetic source identity is introduced at acquisition and persists through every later stage.

\subsubsection{Non-Functional Requirements}

Security requirements include fail-closed input handling, conservative identifier formats, bounded history, source enumeration, immutable pack definitions, disabled backup, and no sensitive runtime permissions. Privacy requirements include local processing, no remote evidence service, no analytics or advertising SDK, and a clear deletion control. Reliability requirements include deterministic results for identical admitted inputs, replay replacement, pack-specific persistence, and migration from the previous repository schema.

Usability requirements include a clear initial action, a stable three-destination navigation model, non-colour status semantics, concise summaries, optional detail, meaningful accessibility labels, and separation of destructive controls. Maintainability requirements include one authoritative compliance path, typed interfaces, registry-based packs, reusable theme constants, tests for new packs, and removal of duplicate legacy services. Deployability requirements include a stable application identifier, explicit version, current target SDK, release APK and AAB generation, upload-key configuration, privacy-policy content, and documented owner actions.

Performance is treated as a constraint rather than a validated claim. The decision engine operates on bounded in-memory arrays and simple rule lookup. Storage is local and capped per package-permission key. WorkManager is selected for deferrable periodic work rather than exact scheduling. No quantitative battery, long-duration memory, or startup objective is claimed in the final delivery because the available evidence does not cover those dimensions broadly enough.

\subsection{Evidence and Regulation-Pack Design}

This section specifies the two contracts at the centre of the framework. The evidence contract defines what the engine may know; the regulation-pack contract defines which reviewed rule data may be applied. Admission, temporal isolation, persistence, and abuse controls preserve the boundary between them.

Table~\ref{tab:legal-traceability} prevents the pack's legal references from being mistaken for derivation of its numerical thresholds. The cited provisions motivate review questions. They do not contain daily or per-minute permission limits, and the application cannot infer the legally decisive facts listed in the fourth column.

\begin{table}[H]
\centering
\caption{EU GDPR pack: legal motivation and prohibited inference}
\label{tab:legal-traceability}
\scriptsize
\begin{tabular}{>{\raggedright\arraybackslash}p{0.12\textwidth}>{\raggedright\arraybackslash}p{0.15\textwidth}>{\raggedright\arraybackslash}p{0.19\textwidth}>{\raggedright\arraybackslash}p{0.23\textwidth}>{\raggedright\arraybackslash}p{0.15\textwidth}}
\toprule
Legal anchor & Observed input & Permitted engineering prompt & Facts still required & Prohibited conclusion \\
\midrule
Art. 5(1)(b), purpose limitation & Contacts count and time window & Ask whether access serves an explicit and compatible purpose & Purpose specification, compatibility analysis, recipients, user expectation & Contacts access breached purpose limitation \\
Art. 5(1)(c), data minimisation & Location, microphone, or contacts frequency & Ask whether the observed pattern was necessary and proportionate to the stated function & Function, alternatives, precision, duration, retention, data actually produced & A threshold exceedance was unlawful or excessive \\
Art. 6 and, where applicable, Art. 9 & Declared lawful-basis and special-category context fields & Expose a missing basis or Article 9 condition for human review & Controller assessment, validity of basis, data category, exceptions, withdrawal and objection facts & Processing lacked a valid legal basis \\
Art. 5(2) and Art. 25 & Source, pack version, rationale, missing fields, local controls & Preserve a reproducible prompt and evidence boundary & Wider technical and organisational measures, risk, state of the art, records, governance & The controller demonstrated compliance or the App is certified \\
\bottomrule
\end{tabular}
\end{table}

The legal traceability matrix is deliberately asymmetric. Technical evidence may trigger a question, but a legal proposition is not allowed to flow backwards into the evidence record as if it had been observed. This follows the GDPR and final EDPB guidance while retaining the privacy-engineering distinction between a data action, its context, and the problems it may create \cite{eurlex2016,edpb2019,nist8062}.

\subsubsection{Evidence Schema and Provenance}

The evidence schema is designed around reconstructability. A permission summary includes a package identifier, supported permission key, non-negative safe-integer count, start and end of the aggregate window, source enumeration, and optional processing context. The source is one of the controlled values supported by the application, such as native bridge, imported evidence, or simulator. A free-form source string is rejected because it could allow synthetic data to impersonate device evidence.

Context fields represent what a legal review would need but a permission count cannot provide. Examples include purpose, controller identity, lawful basis, special-category condition, foreground relationship, consent state, retention statement, and sharing information. Their absence does not make the technical record invalid. It changes the finding by adding missing-context questions and preventing a legal conclusion. This distinction separates data-quality failure from legal incompleteness.

A finding contains the admitted evidence plus decision metadata: status, risk, rule identifier, threshold, rationale, regulation-pack identifier and version, official source, evaluation time, and missing-context fields. The schema is intentionally verbose. Storage efficiency is less important than the ability to explain why a result appeared and which assumptions applied.

Provenance also constrains metrics. Simulator precision and recall are computed only for records whose expected labels are part of the controlled generator. Imported or native records do not receive an accuracy label because their legal and behavioural ground truth is unknown. Aggregating these sources into one accuracy metric would create a number with no coherent denominator.

\subsubsection{Admission and Fail-Closed Validation}

Validation occurs at runtime even though the application is written in TypeScript. Static types protect code paths compiled from known sources, but values crossing a native bridge, storage migration, import boundary, or untyped caller can violate those types. The validator therefore treats each field as unknown until checked.

Counts and timestamps must be finite JavaScript safe integers. This excludes negative values, fractional event counts, infinities, not-a-number values, and integers whose arithmetic cannot be represented exactly. The time window must have a strictly earlier start and a later end. An evaluation timestamp must be consistent with the window. Retention values must be positive and bounded. Package identifiers and permission keys must match conservative patterns rather than merely being non-empty strings.

Context validation follows the same rule. String fields must actually be strings before trimming; booleans must not be inferred from truthy text or numbers. The engine does not coerce a value such as the text false into a Boolean because this can reverse meaning. Unknown source and permission values produce an invalid or review-hold state instead of falling through to a default threshold.

Fail-closed does not mean the application crashes. It means the system returns a safe, explicit result that prevents incomplete evidence from being described as no concern. The UI can then explain which field was rejected or why review cannot proceed. This approach supports both security and diagnosability.

\subsubsection{Temporal Isolation Model}

Aggregate windows create a double-counting problem. Suppose two records report counts for intervals that overlap. Without event identifiers, the system cannot know whether events in the shared interval appear in both counts. Summing them could manufacture a cross-window threshold. the final delivery therefore aggregates only completed, non-overlapping windows for the same package and permission.

For windows \(w_i=[s_i,e_i)\) and \(w_j=[s_j,e_j)\), compatibility is defined as

\begin{equation}
e_i \leq s_j \quad \lor \quad e_j \leq s_i.
\end{equation}

Adjacency is allowed because half-open intervals that meet at one boundary do not overlap. Exact replay is detected by equal package, permission, start, and end. The newer replay replaces the stored value rather than increasing the history size or total. This makes processing idempotent for a repeated aggregate export.

Retention uses a monotonic watermark derived from the latest admitted window end for the key. Basing expiry on a newly arrived record's timestamp alone would allow a late or forged old record to move the horizon backwards. Records older than the bounded horizon relative to the watermark are removed. Each package-permission history is also capped at a finite maximum. The cap is a denial-of-service safeguard, not a claim about the maximum legitimate event rate.

Cross-window detection is consequently conservative. It may fail to combine useful but overlapping aggregates, yet it avoids asserting a total that the evidence does not support. Event-level identifiers or disjoint platform buckets would permit a richer model. The current design states the limitation rather than guessing how much overlap occurred.

\subsubsection{Rule-Pack Contract in Detail}

The common pack contract separates descriptive metadata from executable rule data. Descriptive metadata includes identifier, display name, jurisdiction label, version, official source, effective or review date, and caveat. Rule data binds a supported permission signal to a research baseline, threshold logic, severity, rationale, reference label, and missing-context questions.

The registry is the only route from a persisted identifier to a pack object. This prevents arbitrary downloaded or user-entered text from being executed as a rule. Pack objects are immutable so an evaluation cannot observe a threshold changing midway through a run. The engine receives the pack as an explicit dependency, which makes test fixtures clear and avoids hidden global state.

Pack identity is copied into findings rather than represented only by a current settings value. If a user evaluates under the EU pack, selects the Research Baseline, and later opens an older finding, the finding still displays the EU pack and version. This historical stability is essential for auditability.

The pack contract deliberately does not require every jurisdiction to supply an equivalent article number. A rule can use a provision, principle, official guidance reference, or non-legal research identifier. The UI renders a general reference label and source link. This avoids encoding the assumption that all legal systems share the GDPR's structure.

\subsubsection{State, Persistence, and Migration}

The privacy context coordinates bridge capability, active pack, loading state, findings, demonstration status, and repository actions. State transitions are explicit: idle, reviewing, completed, and failed or unavailable. The interface disables duplicate actions while a review is running and provides a human-readable result when acquisition returns no evidence.

AsyncStorage is used for app-private persistence. The repository key is versioned. On load, older records are normalised into the current schema where safe. Fields that cannot be reconstructed are marked missing rather than invented. A migration should preserve source identity and original timestamps. If a record cannot satisfy current admission rules, it should be excluded from the active projection and surfaced as a migration issue where appropriate.

Clearing findings removes review records while preserving the application itself. Pack selection may be retained separately because it is a preference, but the UI states what the clear action affects. Android backup is disabled so uninstalling or clearing application data has the expected local-deletion meaning within the platform boundary.

\subsubsection{Security and Abuse Cases}

The first abuse case is source forgery: a caller marks synthetic input as native. Enumeration and bridge-controlled construction reduce this risk inside the application, although a modified binary is outside the threat model. The second is numeric abuse: extremely large or non-finite counts cause overflow or misleading totals. Safe-integer validation rejects them. The third is replay amplification. Exact-window replacement provides idempotence. The fourth is overlap amplification, addressed by temporal isolation.

The fifth case is history exhaustion. An attacker or faulty source supplies many distinct windows. Per-key caps and retention limit memory growth. The sixth is identifier pollution, where many malformed package or permission strings create unbounded partitions. Conservative formats and supported permission enumeration reject these values. The seventh is pack confusion, where an unknown identifier is presented under a trusted legal name. Registry-only loading and copied pack metadata prevent silent fallback.

The final abuse case is epistemic rather than computational: a user or stakeholder treats a finding as a legal determination. Interface caveats, neutral statuses, official-source links, and missing-context questions reduce this risk. They cannot eliminate it. Training, documentation, and qualified review remain necessary controls.

\subsection{Traceability and Evaluation Design}

The architecture is complete only when its requirements can be tested at the appropriate layer. This final section links requirements to evidence, identifies stable invariants, and records how future extensions should renew acceptance.

\subsubsection{Traceability from Requirement to Evidence}

Every requirement is paired with a planned evidence source. Deterministic evaluation is covered by seeded and exact-boundary fixtures. Provenance preservation is covered by type, repository, and presentation tests. Unsafe inputs are covered by adversarial runtime cases. Pack switching is covered by registry and identity assertions plus Settings rendering. Synthetic separation is covered by simulator-source tests and populated-device screenshots.

Navigation and visual hierarchy are evidenced by UI trees and physical screenshots, while accessibility claims remain limited to source and rendered cues. Local processing and permission minimisation are evidenced by dependency review, merged manifest, and installed-package inspection. Delivery is evidenced by successful release tasks, artefact hashes, package metadata, and installation. The matrix does not list legal compliance or broad user adoption because no test in the current protocol can establish them.

\subsubsection{Evaluation Strategy}

Evaluation is deliberately layered. Source-level checks establish type and lint conformance. Deterministic tests establish consistency with the published rule specification and adversarial rejection behaviour. Android builds establish package integration. Manifest and package inspection establish requested-permission facts. ADB interaction and screenshots establish bounded physical rendering and launch behaviour. Each result is reported at its own layer; no result is extrapolated into legal compliance or population effectiveness.

\subsubsection{Design Invariants}

The architecture is governed by invariants that should remain true even when screens, acquisition connectors, or legal packs change. First, no accepted finding exists without an admitted evidence source and bounded temporal scope. Second, no legal or research rule exists without a pack identifier and version. Third, unknown, malformed, or unavailable evidence cannot be translated into a reassuring status. Fourth, simulator provenance cannot be upgraded by persistence or presentation. Fifth, historical findings retain the rule identity used at evaluation time. Sixth, local deletion remains available without requiring an account or remote approval.

These invariants are more stable than component names. A future implementation could replace AsyncStorage with Room, React Native with another presentation framework, or bundled packs with signed declarative packages while preserving the same assurance properties. Conversely, a refactor that keeps every current class name but drops source identity or silent-fallback protection would violate the design. Stating invariants therefore makes architectural review less dependent on the incidental code structure of the final delivery.

The invariants also guide test selection. Admission tests challenge the first and third invariants. Pack-switch and persistence tests challenge the second and fifth. Simulator and cross-window tests challenge the fourth. Manifest, backup, and clear-data checks contribute to the sixth. This mapping helps maintainers decide which evidence to renew after a change. It also reveals gaps: authentication of an external producer, for example, is not yet an invariant guaranteed by the prototype and must not be implied by the controlled source label.

An invariant can fail through wording as well as computation. If a finding correctly stores evidence gap but the screen headlines the review as safe, the third invariant has failed at presentation. If an export omits the pack version, the fifth has failed outside the repository. End-to-end checks are therefore necessary even when unit tests are comprehensive.

Future release reviews should record each invariant as satisfied, failed, or not evidenced, with direct artefact references. A failed invariant blocks promotion. A not-evidenced item may permit a narrower research build, but its associated claim must remain absent. This practice keeps the product's trust model explicit as functionality grows.

\subsubsection{Decision Records for Future Extension}

Future contributors should document material changes as concise decision records. A record should state the problem, evidence, alternatives, selected option, affected invariants, privacy consequences, pack consequences, migration, rollback, and verification required. This is especially important for changes that appear convenient but alter authority, such as adding remote telemetry, inferring jurisdiction automatically, importing an unreviewed rules file, or expanding Android privileges.

A decision record does not replace code review or legal review. It makes assumptions visible before they become distributed across implementation and copy. It also provides a durable explanation when later maintainers ask why the product refuses a seemingly simple shortcut. In a thesis prototype intended for store progression, this continuity matters because deployment work will occur after the original evaluation snapshot.

The first required record for a new acquisition connector should describe deployment authority and observed coverage. The first record for a new regulation pack should identify responsible reviewers and effective sources. The first record for remote services should include a data-flow and retention analysis. The first record for a new user-facing status should specify evidence semantics and comprehension tests. These prompts keep expansion aligned with the central evidence-bounded principle.

\subsection{Temporal-combination model}

The final delivery adds a narrow temporal inference layer between observation capture and the human-review workflow. Let an observation be $e=(p,t,k,c,s)$, where $p$ is a package-scoped subject, $t$ an observed time, $k$ a normalised observation type, $c$ a non-negative count, and $s$ an evidence source. The domain of $k$ includes permission-relevant data and sensor activities such as location, microphone, contacts, camera, body sensors, activity recognition, media-related access, clipboard reads, device identifiers, backgrounding, and data transfer. It is an engineering vocabulary, not a claim that every Android API call has been observed or legally classified.

A temporal profile $r=(id,W,M,L,q)$ is supplied by the active regulation pack, not by the generic detector. $W$ is the rule-specific positive window length, $M$ is a multiset of required types with minimum counts, $L$ contains the rule-owned legal references, and $q$ is a short rationale. For a candidate time $T$, the detector considers the per-package ledger $E_{p,T,W}=\{e \mid T-W \leq t \leq T\}$. A profile matches precisely when, for every $(k,n)\in M$, the summed count of observations of type $k$ in $E_{p,T,W}$ is at least $n$. No requirement is imposed on the order in which the members of $M$ occur. This makes a combination such as location, activity recognition, and body-sensor access equivalent under every permutation while retaining its observed timestamps and sources as evidence.

The formulation avoids a global fixed time threshold. $W$, $M$, $L$, and $q$ are properties of the particular assessment item. The mounted GDPR pack expresses distinct wearable location--health, multimodal biometric-capture, cross-domain-profiling, post-background-media-access, and high-frequency-location profiles; their durations and required observations differ. The profile is only an alerting hypothesis with cited legal context. It is not a conclusion that the processing is unlawful, sensitive in fact, or intended for profiling.

Two boundaries are essential. A match is scoped to one package and a bounded ledger, so an event from one application cannot complete a combination for another. A notification does not change the App's access-control outcome. It appends an evidence-bearing finding for human review and retains non-blocking behaviour. These boundaries preserve the distinction between a technical signal, missing contextual facts, qualified legal assessment, and any eventual human action.

\subsection{Restart and controlled-fixture contracts}

The ledger snapshot contains only package name, pack identifier and version, normalised type, timestamp, count, source category, and a deterministic deduplication key. At restoration the engine rejects the whole snapshot when its schema, pack identity, version, type, time, count, source, or deduplication key is invalid. It also discards events outside the current pack's maximum window or later than the restart time. Clearing local findings clears this temporal state. A controlled fixture is generated from the compiled active profile, not from a hard-coded GDPR combination, and must carry both the controlled-demo label and native-bridge source. The release bridge rejects this test route.
